% 实验报告起步模板 —— 复制到本次实验的工作目录后填写
% 编译：python <skill>/scripts/build.py report.tex --thumbs preview
% 规则：所有测量数据来自 data/*.csv（照片逐格转写、用户核对过）；正文数字用 results.tex 里的宏，不手抄。
\documentclass{labreport}        % 不要外框：\documentclass[plain]{labreport}
\InputIfFileExists{results.tex}{}{}

\labset{
  school     = {XX大学},           % 封面大标题 = school + title
  cover      = {simple},           % simple：通用封面；szu：深圳大学样式
  coursename = {大学物理实验},
  coursecode = {},
  expname    = {实验名称},
  college    = {},
  teacher    = {},
  author     = {},
  studentid  = {},
  group      = {},
  location   = {},
  date       = {},                 % 不知道就留空，不要编
  submitdate = {},
}

\begin{document}
\makecover

\begin{labsection}{实验目的}
1．……\par
2．……
\end{labsection}

\begin{labsection}{实验原理}
\labsub{……}
……
\begin{equation}
  y = kx + b \label{eq:main}
\end{equation}
式中 ……（每个符号的含义和单位）。
\end{labsection}

\begin{labsection}{实验仪器}
% 仪器表：名称 / 型号 / 量程 / 分度值 / 仪器误差限 Δ仪（查不到的项问用户，不编）
% \input{tab/instr.tex}       ← csv2tex.py 生成
\end{labsection}

\begin{labsection}{实验内容与步骤}
1．……\par
\textbf{注意事项：}……
\end{labsection}

\begin{labsection}{数据记录和处理}
\textbf{1．数据记录}\par
% \input{tab/t1.tex}

\textbf{2．数据处理}\par
% 代入数字的计算示例 → 图 → 拟合 / 逐差结果 → 不确定度 → 最终结果
% \labfig[0.7\linewidth]{fig/fig1.pdf}{$y$-$x$ 关系曲线}[fig:1]
% 拟合得 $k=\num{\FitK}$，$r=\num{\FitR}$。      ← fit_plot.py --macros results.tex --name Fit
% 结果：$x=(\XVal\pm\XU)$\,单位（$P=0.683$），相对不确定度 \XRel\%。   ← uncert.py --macros results.tex --name X
……
\end{labsection}

\begin{labsection}{结果陈述及总结}
1．由表 \ref{tab:1} 及图 \ref{fig:1} 可知……（每条结论带数据：线性给 $r$，比较给相对误差）\par
2．误差分析：哪个操作 → 影响哪个量 → 偏大还是偏小 → 怎么改进。
\end{labsection}

\begin{labsection}{思考题}
\question{题目原文}
……
\end{labsection}

\labcomment
\labgrading{预习/20,操作及记录/40,数据处理与结果陈述/30,思考题/10}

% 原始数据照片、预习试卷：先 python <skill>/scripts/prep_photos.py ... --out raw
% \labphotos{原始数据}{raw/raw1.jpg,raw/raw2.jpg}
% \labphotos{预习试卷}{raw/prep1.jpg}
\end{document}
